\textbf{Phần 1. Trắc nghiệm lựa chọn}
\setcounterpageref{ex}{0}

\begin{ex} %[4D3N15]
Dãy số nào sau đây là dãy số tự nhiên?
\choice
{$1; 2; 3; 4; 5; \dots$}
{\True $0; 1; 2; 3; 4; 5; \dots$}
{$0; 2; 4; 6; 8; \dots$}
{$1; 3; 5; 7; 9; \dots$}
\loigiai{
Dãy số tự nhiên bắt đầu từ số 0: $0; 1; 2; 3; 4; 5; \dots$\\
\textbf{Đáp án của câu hỏi là B}
}
\end{ex}

\begin{ex} %[4D3N15]
Số tự nhiên nhỏ nhất là
\choice
{\True $0$}
{$1$}
{Không có}
{$10$}
\loigiai{
Số tự nhiên nhỏ nhất là số 0.\\
\textbf{Đáp án của câu hỏi là A}
}
\end{ex}

\begin{ex} %[4D3H15]
Có bao nhiêu số tự nhiên lớn nhất?
\choice
{Chỉ có số $999\,999\,999$}
{\True Không có số tự nhiên lớn nhất}
{Chỉ có số 1 tỉ}
{Có 1 số}
\loigiai{
Dãy số tự nhiên kéo dài vô tận, không có số tự nhiên lớn nhất.\\
\textbf{Đáp án của câu hỏi là B}
}
\end{ex}

\begin{ex} %[4D3H15]
Hai số tự nhiên liên tiếp hơn kém nhau bao nhiêu đơn vị?
\choice
{$0$ đơn vị}
{$2$ đơn vị}
{$10$ đơn vị}
{\True $1$ đơn vị}
\loigiai{
Hai số tự nhiên liên tiếp hơn kém nhau 1 đơn vị.\\
\textbf{Đáp án của câu hỏi là D}
}
\end{ex}

\begin{ex} %[4D3H15]
Số liền sau của số $n$ (với $n$ là số tự nhiên) là
\choice
{$n - 1$}
{$n \times 2$}
{\True $n + 1$}
{$n : 2$}
\loigiai{
Số liền sau của $n$ là $n + 1$.\\
\textbf{Đáp án của câu hỏi là C}
}
\end{ex}

\begin{ex} %[4D3H15]
Dãy số tự nhiên chẵn liên tiếp bắt đầu từ 0 là
\choice
{$1; 3; 5; 7; \dots$}
{\True $0; 2; 4; 6; 8; \dots$}
{$2; 4; 6; 8; 10$}
{$0; 1; 2; 3; \dots$}
\loigiai{
Dãy số chẵn: $0; 2; 4; 6; 8; \dots$\\
\textbf{Đáp án của câu hỏi là B}
}
\end{ex}

\begin{ex} %[4D3V15]
Cho dãy số: $2; 5; 8; 11; 14; \dots$ Quy luật của dãy số này là
\choice
{Mỗi số gấp đôi số đứng trước}
{Mỗi số bằng số liền trước cộng thêm 2}
{Mỗi số bằng số liền trước nhân với 3}
{\True Mỗi số hạng (kể từ số thứ hai) bằng số liền trước cộng thêm 3}
\loigiai{
Khoảng cách giữa hai số liên tiếp là 3 đơn vị.\\
\textbf{Đáp án của câu hỏi là D}
}
\end{ex}

\begin{ex} %[4D3V15]
Tìm số hạng thứ 6 của dãy số: $1; 4; 7; 10; 13; \dots$
\choice
{\True $16$}
{$15$}
{$17$}
{$18$}
\loigiai{
Số thứ 6 là $13 + 3 = 16$.\\
\textbf{Đáp án của câu hỏi là A}
}
\end{ex}

\begin{ex} %[4D3V15]
Dãy số $10; 12; 14; \dots; 98$ có tất cả bao nhiêu số hạng?
\choice
{$44$ số}
{$46$ số}
{$88$ số}
{\True $45$ số}
\loigiai{
Số số hạng: $(98 - 10) : 2 + 1 = 45$ số hạng.\\
\textbf{Đáp án của câu hỏi là D}
}
\end{ex}

\begin{ex} %[4D3V15]
Tính tổng của dãy số: $1 + 2 + 3 + \dots + 20$.
\choice
{$200$}
{$220$}
{\True $210$}
{$190$}
\loigiai{
$(1 + 20) \times 20 : 2 = 210$.\\
\textbf{Đáp án của câu hỏi là C}
}
\end{ex}

\begin{ex} %[4D3V15]
Tìm số thích hợp điền vào dấu hỏi chấm: $1; 2; 4; 8; 16; ?$
\choice
{\True $32$}
{$24$}
{$30$}
{$64$}
\loigiai{
Mỗi số gấp đôi số đứng trước nó: $16 \times 2 = 32$.\\
\textbf{Đáp án của câu hỏi là A}
}
\end{ex}

\begin{ex} %[4D3C15]
Để đánh số trang một cuốn sách từ trang 1 đến trang 50 người ta cần dùng tất cả bao nhiêu chữ số?
\choice
{$90$ chữ số}
{$95$ chữ số}
{\True $91$ chữ số}
{$100$ chữ số}
\loigiai{
Từ trang 1 đến 9: 9 chữ số.\\
Từ trang 10 đến 50: $(50 - 10 + 1) \times 2 = 41 \times 2 = 82$ chữ số.\\
Tổng số chữ số: $9 + 82 = 91$ chữ số.\\
\textbf{Đáp án của câu hỏi là C}
}
\end{ex}

\textbf{Phần 2. Trắc nghiệm đúng sai}
\setcounterpageref{ex}{0}

\begin{ex} %[4D3H15]
Xét tính đúng sai về các đặc điểm của dãy số tự nhiên:
\choiceTF
{\True Số 0 là số tự nhiên nhỏ nhất.}
{Số 1 là số tự nhiên nhỏ nhất.}
{\True Dãy số tự nhiên có vô số phần tử (không có số lớn nhất).}
{\True Giữa hai số tự nhiên liên tiếp không có số tự nhiên nào khác.}
\loigiai{
\begin{itemize}
    \item 0 là số tự nhiên nhỏ nhất. Đúng.
    \item 1 không phải số tự nhiên nhỏ nhất (0 nhỏ hơn). Sai.
    \item Không có số lớn nhất: Đúng.
    \item Hai số tự nhiên liên tiếp liền nhau: Đúng.
\end{itemize}
\begin{itemize}
    \item a) \textbf{Đúng}.
    \item b) \textbf{Sai}.
    \item c) \textbf{Đúng}.
    \item d) \textbf{Đúng}.
\end{itemize}
}
\end{ex}

\begin{ex} %[4D3V15]
Cho dãy số: $3; 6; 9; 12; \dots; 60$. Xét tính đúng sai:
\choiceTF
{\True Các số trong dãy đều chia hết cho 3.}
{\True Dãy số có tất cả $20$ số hạng.}
{Tổng của số đầu và số cuối là $60$.}
{\True Tổng của tất cả các số hạng trong dãy là $630$.}
\loigiai{
\begin{itemize}
    \item Các số chia hết cho 3: Đúng.
    \item Số số hạng: $(60 - 3) : 3 + 1 = 20$. Đúng.
    \item Số đầu $+$ Số cuối $= 3 + 60 = 63 \neq 60$. Sai.
    \item Tổng $= 63 \times 20 : 2 = 630$. Đúng.
\end{itemize}
\begin{itemize}
    \item a) \textbf{Đúng}.
    \item b) \textbf{Đúng}.
    \item c) \textbf{Sai}.
    \item d) \textbf{Đúng}.
\end{itemize}
}
\end{ex}

\textbf{Phần 3. Trả lời ngắn}
\setcounterpageref{ex}{0}

\begin{ex} %[4D3H15]
Tìm số liền trước của số tự nhiên $a$ với $a = 100$.
\shortans{99}
\loigiai{
\textbf{Đáp số:} $99$\\
$100 - 1 = 99$.
}
\end{ex}

\begin{ex} %[4D3H15]
Tìm số hạng tiếp theo của dãy số: $5; 10; 15; 20; \dots$
\shortans{25}
\loigiai{
\textbf{Đáp số:} $25$\\
$20 + 5 = 25$.
}
\end{ex}

\begin{ex} %[4D3V15]
Có bao nhiêu số tự nhiên có 3 chữ số?
\shortans{900}
\loigiai{
\textbf{Đáp số:} $900$\\
$999 - 100 + 1 = 900$.
}
\end{ex}

\begin{ex} %[4D3V15]
Tính tổng của 10 số tự nhiên liên tiếp đầu tiên ($0; 1; 2; \dots; 9$).
\shortans{45}
\loigiai{
\textbf{Đáp số:} $45$\\
$(0 + 9) \times 10 : 2 = 45$.
}
\end{ex}

\begin{ex} %[4D3V15]
Tìm số hạng thứ 10 của dãy số: $2; 4; 6; 8; \dots$
\shortans{20}
\loigiai{
\textbf{Đáp số:} $20$\\
$2 \times 10 = 20$.
}
\end{ex}

\begin{ex} %[4D3C15]
Để viết các số tự nhiên liên tiếp từ 1 đến 100 cần dùng tất cả bao nhiêu chữ số?
\shortans{192}
\loigiai{
\textbf{Đáp số:} $192$\\
$1 \rightarrow 9$: 9 chữ số.\\
$10 \rightarrow 99$: $90 \times 2 = 180$ chữ số.\\
$100$: 3 chữ số.\\
Tổng: $9 + 180 + 3 = 192$ chữ số.
}
\end{ex}
